\Askhsh{29837_SOLUTION}{1}
\begin{ylh}{1}
    \item [Ύλη από εικόνα]
\end{ylh}
ΛΥΣΗ
\begin{alist}
\item [α)] Έχουμε την $f(x) = \frac{1}{1-x}$ με $D_f = \mathbb{R}-\{1\}$. Έστω $x_1, x_2 \in D_f$ ώστε:
\[ f(x_1) = f(x_2) \Leftrightarrow \frac{1}{1-x_1} = \frac{1}{1-x_2} \Leftrightarrow 1-x_2 = 1-x_1 \Leftrightarrow x_1 = x_2, \]
οπότε η συνάρτηση $f$ είναι '1-1' επομένως αντιστρέφεται.

Για την εύρεση του τύπου της αντιστρόφου γνωρίζουμε ότι $y = f(x) \Leftrightarrow f^{-1}(y) = x$.
Οπότε για $x \in D_f$, θέτουμε $y = f(x)$ και λύνουμε ως προς $x$, παίρνοντας τους περιορισμούς για το $y$, που θα ορίσουν το σύνολο τιμών της $f$, που αποτελεί το πεδίο ορισμού της αντιστρόφου. Είναι για $x \ne 1$ (1):
\[ y = f(x) \Leftrightarrow y = \frac{1}{1-x} \Leftrightarrow 1-x = \frac{1}{y}, y \ne 0 \Leftrightarrow x = 1-\frac{1}{y}, y \ne 0,\text{που ικανοποιεί την (1)} \Leftrightarrow x = \frac{y-1}{y}, y \ne 0. \]
Επομένως $f^{-1}(y) = \frac{y-1}{y}, y \ne 0$ ή $f^{-1}(x) = \frac{x-1}{x}, x \ne 0$.

\item [β)] Για την $f \circ f$ είναι: $A_1 = \{x \in D_f : f(x) \in D_f\} = \{x \in \mathbb{R}-\{1\} : f(x) \in \mathbb{R}-\{1\}\}$
\[ = \left\{x \ne 1 : \frac{1}{1-x} \ne 1\right\} = \{x \ne 1 : 1 \ne 1-x\} = \{x \ne 1 : x \ne 0\} = \mathbb{R}-\{0,1\} \ne \emptyset, \]
οπότε ορίζεται η σύνθεση $f \circ f$
με τύπο: $(f \circ f)(x) = f(f(x)) = \frac{1}{1-f(x)} = \frac{1}{1-\frac{1}{1-x}} = \frac{1}{\frac{1-x-1}{1-x}} = \frac{1-x}{-x} = \frac{x-1}{x}$ με $x \in \mathbb{R}-\{0,1\}$.

\item [γ)] Λόγω των ερωτημάτων (α), (β) έχουμε:
$f^{-1}(x) = \frac{x-1}{x}$ με $D_{f^{-1}} = \mathbb{R}^* = (-\infty,0) \cup (0,+\infty)$ και
$(f \circ f)(x) = \frac{x-1}{x}$ με $D_{f \circ f} = \mathbb{R}-\{0,1\} = (-\infty,0) \cup (0,1) \cup (1,+\infty)$.

Εφόσον $D_{f^{-1}} \ne D_{f \circ f}$, οι συναρτήσεις $f \circ f$ και $f^{-1}$ δεν είναι ίσες. Επομένως δεν συμφωνούμε με τον ισχυρισμό του μαθητή.

\item [δ)] Είναι:
\[ \int_2^3 \varphi(x)dx = \int_2^3 \frac{x-1}{x}dx = \int_2^3 \left(1-\frac{1}{x}\right)dx = [x-\ln|x|]_2^3 = (3-\ln 3)-(2-\ln 2) = 1+\ln 2-\ln 3. \]
\end{alist}